\section{Introduction}
\label{sec:introduction}

Network slicing is a central mechanism for supporting diverse service requirements in 5G and emerging 6G systems. A slice may be optimized for enhanced mobile broadband (eMBB), ultra-reliable low-latency communication (URLLC), massive IoT, or an operator-defined service class. In practice, however, a slice is not configured at a single point. The radio access network (RAN) must allocate radio resources and scheduling priority, the core network must bind subscribers and PDU sessions to QoS flows, and the cloud-native infrastructure must deploy and update network functions. This cross-domain nature creates a persistent gap between service-level objectives and the low-level knobs exposed by individual systems.

Intent-driven management addresses this gap by allowing a consumer to state the desired outcome rather than the sequence of operations required to achieve it. 3GPP TS 28.312 defines intent as an outcome-oriented management construct and discusses the lifecycle of intent handling, including translation, fulfillment, reporting, and conflict handling~\cite{3gpp_ts28312}. O-RAN similarly introduces RAN automation functions such as rApps for policy and optimization~\cite{oran_oad_r005}. Yet an open-source, end-to-end implementation that connects high-level intents to both free5GC and OCUDU remains difficult because the relevant controls live across different configuration systems, APIs, and reconciliation loops.

This work builds an intent-driven orchestration system for open-source 5G network slices. The design uses Kubernetes as the intent interface and execution substrate. A consumer submits an \texttt{E2EQoSIntent} custom resource that describes slice expectations, such as \texttt{bandwidth: high}, \texttt{latency: ultra-low}, and \texttt{resourceShare: Full}. The orchestrator translates these expectations into domain parameters: 5QI and QFI for the 5G core, plus S-NSSAI, PRB policy ratios, and scheduling priority for the RAN. RAN changes are handled as declarative Nephio package mutations and realized by the OCUDU Operator. In the core domain, the official Nephio free5GC Operator reconciles network-function deployments, while subscriber registration and QoS parameters are applied through free5GC-facing configuration logic.

The contributions are:
\begin{itemize}
  \item We design a Kubernetes-native intent model, \texttt{E2EQoSIntent}, for describing multiple 5G slice objectives in one end-to-end resource.
  \item We implement a translation layer from service-level expectations to concrete free5GC and OCUDU parameters, including 5QI, QFI, S-NSSAI, PRB policy ratio, and priority.
  \item An OCUDU Operator provides hierarchical CRDs, cross-resource validation, configuration rendering, and RAN reconciliation.
  \item We integrate Nephio/GitOps and domain operators into RAN and core NSSMF paths, with direct free5GC configuration for dynamic subscriber and QoS state.
  \item We evaluate the prototype with eMBB and URLLC slices and report both data-plane differentiation and management-plane orchestration latency.
\end{itemize}
